% GENERATED FILE. Edit canonical lessons, then run npm run generate:books.
% canonical-source-hash: fnv1a64:19639ede38bf3ccb
% canonical-lessons: GU-C243-path, GU-C243-gokhvu, GU-C243-copdi, GU-C243-daphtar, GU-C243-hajar

\chapter{Lessons and school bags}
\label{ch:gu-lessons-and-school-bags}
\hlchaptermodality{243}

\begin{chapteropening}
\textbf{By the end of this chapter:} \emph{I can talk about lessons, schoolbooks and who is present in class.}

Complete the last lesson of chapter 243: I can talk about lessons, schoolbooks and who is present in class.
\end{chapteropening}

\section[\texorpdfstring{\gu{પાઠ} (pāṭh)}{પાઠ (pāṭh)}]{\gu{પાઠ} (pāṭh) — lesson, reading}

\label{lesson:GU-C243-path}



\begin{quote}
Before the new one: say the Gujarati for to believe, then the Gujarati for to change.
\end{quote}

\subsection*{You'll want to know: \gu{પાઠ}}
\textbf{\gu{પાઠ}} — \emph{pāṭh} — \textquotedblleft{}lesson, reading\textquotedblright{}.

\begin{quote}
\textbf{\gu{આજે} \gu{અમે} \gu{નવો} \gu{પાઠ} \gu{ભણીએ} \gu{છીએ}.} — \emph{āje ame navo pāṭh bhaṇīe chīe} — \textquotedblleft{}Today we study a new lesson.\textquotedblright{}
\end{quote}

\begin{cousinweb}
\textbf{\gu{પાઠ}} (\emph{pāṭh}) is Sanskrit \emph{pāṭha}, \textquotedblleft{}text of a book; study\textquotedblright{}, taken over whole as a learned loan, a tatsama (CDIAL 8035).
\end{cousinweb}

\begin{grammarlens}[title={What you've built}]
One more word for everyday life at A2.
\end{grammarlens}

\subsection*{Your turn}
\begin{itemize}
\raggedright
  \item \emph{Say it:} \emph{pāṭh}
  \item \emph{Say it:} \emph{pāṭh}, once more
  \item \emph{Say it:} \emph{āje ame navo pāṭh bhaṇīe chīe} — Today we study a new lesson.
  \item \emph{Recall:} say the Gujarati for to believe, then the Gujarati for to change, then say \emph{pāṭh} again
\end{itemize}

\begin{tcolorbox}[breakable,colback=teal!4,colframe=teal!35!black,title={Before you move on}]
What does \gu{પાઠ} mean? (\textquotedblleft{}Lesson, reading\textquotedblright{}.) Say it once more.
\end{tcolorbox}
\section[\texorpdfstring{\gu{ગોખવું} (gokhvũ)}{ગોખવું (gokhvũ)}]{\gu{ગોખવું} (gokhvũ) — to learn by heart, to cram}

\label{lesson:GU-C243-gokhvu}



\begin{quote}
Before the new one: say the Gujarati for to change, then the Gujarati for lesson.
\end{quote}

\subsection*{You'll want to know: \gu{ગોખવું}}
\textbf{\gu{ગોખવું}} — \emph{gokhvũ} — \textquotedblleft{}to learn by heart, to cram\textquotedblright{}.

\begin{quote}
\textbf{\gu{તે} \gu{રોજ} \gu{પાઠ} \gu{ગોખે} \gu{છે}.} — \emph{te roj pāṭh gokhe che} — \textquotedblleft{}He learns the lesson by heart every day.\textquotedblright{}
\end{quote}

\begin{cousinweb}
The origin of \textbf{\gu{ગોખવું}} (\emph{gokhvũ}) is not settled. CDIAL lists no Gujarati verb of this form with the meaning \textquotedblleft{}to learn by rote\textquotedblright{}, and our other sources do not trace it.
\end{cousinweb}

\begin{grammarlens}[title={What you've built}]
One more word for everyday life at A2.
\end{grammarlens}

\subsection*{Your turn}
\begin{itemize}
\raggedright
  \item \emph{Say it:} \emph{gokhvũ}
  \item \emph{Say it:} \emph{gokhvũ}, once more
  \item \emph{Say it:} \emph{te roj pāṭh gokhe che} — He learns the lesson by heart every day.
  \item \emph{Recall:} say the Gujarati for to change, then the Gujarati for lesson, then say \emph{gokhvũ} again
\end{itemize}

\begin{tcolorbox}[breakable,colback=teal!4,colframe=teal!35!black,title={Before you move on}]
What does \gu{ગોખવું} mean? (\textquotedblleft{}To learn by heart, to cram\textquotedblright{}.) Say it once more.
\end{tcolorbox}
\section[\texorpdfstring{\gu{ચોપડી} (copḍī)}{ચોપડી (copḍī)}]{\gu{ચોપડી} (copḍī) — schoolbook, small book}

\label{lesson:GU-C243-copdi}



\begin{quote}
Before the new one: say the Gujarati for lesson, then the Gujarati for to learn by heart.
\end{quote}

\subsection*{You'll want to know: \gu{ચોપડી}}
\textbf{\gu{ચોપડી}} — \emph{copḍī} — \textquotedblleft{}schoolbook, small book\textquotedblright{}.

\begin{quote}
\textbf{\gu{આ} \gu{મારી} \gu{નવી} \gu{ચોપડી} \gu{છે}.} — \emph{ā mārī navī copḍī che} — \textquotedblleft{}This is my new schoolbook.\textquotedblright{}
\end{quote}

\begin{cousinweb}
The origin of \textbf{\gu{ચોપડી}} (\emph{copḍī}) is not settled. CDIAL has no entry for it, and our other sources do not trace it further.
\end{cousinweb}

\begin{grammarlens}[title={What you've built}]
One more word for everyday life at A2.
\end{grammarlens}

\subsection*{Your turn}
\begin{itemize}
\raggedright
  \item \emph{Say it:} \emph{copḍī}
  \item \emph{Say it:} \emph{copḍī}, once more
  \item \emph{Say it:} \emph{ā mārī navī copḍī che} — This is my new schoolbook.
  \item \emph{Recall:} say the Gujarati for lesson, then the Gujarati for to learn by heart, then say \emph{copḍī} again
\end{itemize}

\begin{tcolorbox}[breakable,colback=teal!4,colframe=teal!35!black,title={Before you move on}]
What does \gu{ચોપડી} mean? (\textquotedblleft{}Schoolbook, small book\textquotedblright{}.) Say it once more.
\end{tcolorbox}
\section[\texorpdfstring{\gu{દફતર} (daphtar)}{દફતર (daphtar)}]{\gu{દફતર} (daphtar) — school bag; office file}

\label{lesson:GU-C243-daphtar}



\begin{quote}
Before the new one: say the Gujarati for to learn by heart, then the Gujarati for schoolbook.
\end{quote}

\subsection*{You'll want to know: \gu{દફતર}}
\textbf{\gu{દફતર}} — \emph{daphtar} — \textquotedblleft{}school bag; office file\textquotedblright{}.

\begin{quote}
\textbf{\gu{તેનું} \gu{દફતર} \gu{ભારે} \gu{છે}.} — \emph{tenũ daphtar bhāre che} — \textquotedblleft{}His school bag is heavy.\textquotedblright{}
\end{quote}

\begin{cousinweb}
\textbf{\gu{દફતર}} (\emph{daphtar}) is a loan from Arabic \emph{daftar}, \textquotedblleft{}register, account book\textquotedblright{}, which reached Gujarati through Persian. Lane records the view that Arabic itself took the word from Persian (Lane).
\end{cousinweb}

\begin{grammarlens}[title={What you've built}]
One more word for everyday life at A2.
\end{grammarlens}

\subsection*{Your turn}
\begin{itemize}
\raggedright
  \item \emph{Say it:} \emph{daphtar}
  \item \emph{Say it:} \emph{daphtar}, once more
  \item \emph{Say it:} \emph{tenũ daphtar bhāre che} — His school bag is heavy.
  \item \emph{Recall:} say the Gujarati for to learn by heart, then the Gujarati for schoolbook, then say \emph{daphtar} again
\end{itemize}

\begin{tcolorbox}[breakable,colback=teal!4,colframe=teal!35!black,title={Before you move on}]
What does \gu{દફતર} mean? (\textquotedblleft{}School bag; office file\textquotedblright{}.) Say it once more.
\end{tcolorbox}
\section[\texorpdfstring{\gu{હાજર} (hājar)}{હાજર (hājar)}]{\gu{હાજર} (hājar) — present, in attendance}

\label{lesson:GU-C243-hajar}



\begin{quote}
Before the new one: say the Gujarati for schoolbook, then the Gujarati for school bag.
\end{quote}

\subsection*{You'll want to know: \gu{હાજર}}
\textbf{\gu{હાજર}} — \emph{hājar} — \textquotedblleft{}present, in attendance\textquotedblright{}.

\begin{quote}
\textbf{\gu{બધા} \gu{વિદ્યાર્થી} \gu{વર્ગમાં} \gu{હાજર} \gu{છે}.} — \emph{badhā vidyārthī vargmā̃ hājar che} — \textquotedblleft{}All the students are present in class.\textquotedblright{}
\end{quote}

\begin{cousinweb}
\textbf{\gu{હાજર}} (\emph{hājar}) is a loan from Arabic \emph{ḥāḍir}, \textquotedblleft{}present, at hand\textquotedblright{}, which reached Gujarati through Persian. Its Arabic root is \emph{ḥ-ḍ-r}, \textquotedblleft{}to be present\textquotedblright{} (Lane).
\end{cousinweb}

\begin{grammarlens}[title={What you've built}]
One more word for everyday life at A2.
\end{grammarlens}

\subsection*{Your turn}
\begin{itemize}
\raggedright
  \item \emph{Say it:} \emph{hājar}
  \item \emph{Say it:} \emph{hājar}, once more
  \item \emph{Say it:} \emph{badhā vidyārthī vargmā̃ hājar che} — All the students are present in class.
  \item \emph{Recall:} say the Gujarati for schoolbook, then the Gujarati for school bag, then say \emph{hājar} again
\end{itemize}

\begin{tcolorbox}[breakable,colback=teal!4,colframe=teal!35!black,title={Before you move on}]
What does \gu{હાજર} mean? (\textquotedblleft{}Present, in attendance\textquotedblright{}.) Say it once more.
\end{tcolorbox}
